\documentclass[10pt,a4paper]{amsart}
\usepackage[margin=19mm]{geometry}
\usepackage{amsmath,amssymb,mathtools}
\usepackage[scr=rsfs,cal=euler]{mathalfa}
\usepackage{xfrac}
\usepackage[unicode,hidelinks,bookmarks=false]{hyperref}
\newcommand{\ie}{i.e.\ }
\newcommand{\cf}{cf.\ }
\newcommand{\resp}{respectively\ }
\newcommand{\Subsection}{\subsection}
\providecommand{\nobreakdash}{\nobreak\hbox{-}}
\setlength{\emergencystretch}{2em}
\multlinegap0pt
\usepackage{bm}
\usepackage{bbm}
\usepackage{dsfont}
\usepackage{xspace}
\usepackage{cancel}
\usepackage{mathtools}
\usepackage{enumitem}
\usepackage{amsthm}
\makeatletter
\@ifundefined{theorem}{\newtheorem{theorem}{Theorem}}{}
\@ifundefined{lemma}{\newtheorem{lemma}[theorem]{Lemma}}{}
\makeatother
\newcommand{\sgn}{\mathrm{sgn}}
\title[Highest weight vectors of tensors]{Highest weight vectors of tensors}
\author{Alimzhan Amanov, Damir Yeliussizov}
\date{Source: arXiv:2504.15413v2 (CC BY 4.0)}
\begin{document}
\maketitle

\noindent\emph{Source and licence.} Highest weight vectors of tensors. Version arXiv:2504.15413v2 (CC BY 4.0), distributed under the Creative Commons Attribution 4.0 International licence, as recorded in the arXiv licence metadata for this exact version and shown on the abstract page https://arxiv.org/abs/2504.15413v2. Full rights notice: licenses/arxiv-2504.15413-CC-BY-4.0.txt.\par\smallskip

\noindent\textit{Editorial scope.} Counted unit: the Lemma whose source label is \texttt{lemma:gst} in the article (source0.tex lines 1533--1544, source label \texttt{lemma:gst}), delivered together with the article's own complete original proof of it (source0.tex lines 1546--1638, one \texttt{proof} environment of 93 lines). No mathematics is written, repaired or re-derived: statement and proof are the article's own text line for line. Required context retained uncounted: none. Every definition, display and label of those lines is retained unchanged. Omitted from this package and not redistributed: 1--1532, 1545--1545, 1639--2547. Nothing inside a retained excerpt is omitted or altered: every retained body is delivered exactly as the article's lines are written, so no omission receipt of any kind accompanies this package. Citations: the retained excerpts carry no citation command at all, so no bibliography entry of the article is transcribed and no reference is redistributed. Reference adaptation: none was needed and none was made; every reference inside the delivered unit resolves inside it, so no unresolved reference or citation is produced. Printed slips kept literal: no printed word, symbol, hypothesis or number of the counted statement or of its proof was altered. Layout and class: the article's own class is \texttt{amsart} with packages including hyperref, amsthm, amsfonts, amsmath, amssymb, amscd, verbatim, float, wrapfig, mathtools, todonotes, enumitem, ytableau, tikz; the wrapper is the lane's pinned \texttt{amsart} editorial wrapper at 10pt on A4, which restates the theorem-like environments the retained text needs and adds the article's own macro definitions that the retained text needs (\texttt{\textbackslash sgn}), with extra wrapper packages (bm, bbm, dsfont, xspace, cancel, mathtools, enumitem). The delivered numbering is the unit's own. Assets: no retained span contains a figure environment, a graphics-inclusion command, a PDF-page inclusion, a table or a TikZ picture; no image file is copied into this package, the package declares no asset dependency and no image file is redistributed with it. Licence: version arXiv:2504.15413v2 of this article is distributed under the Creative Commons Attribution 4.0 International licence, as recorded in the arXiv licence metadata for this exact version and shown on the abstract page https://arxiv.org/abs/2504.15413v2. A full rights notice is delivered at licenses/arxiv-2504.15413-CC-BY-4.0.txt. Characters: every retained excerpt is pure ASCII, so the delivered engine, XeLaTeX with its default Latin Modern text font, reports no missing character.
\begin{lemma}\label{lemma:gst}
    Let $\lambda\vdash m$ be a partition of length $\ell = \ell(\lambda)$. For any pair of words $(s,t) \in A(\lambda) \times A(\lambda')$ the following properties hold:
    \begin{enumerate}[label=(\alph*)]
        \item  $G(s,t) \in Y_{s} \backslash S_m / Y_{t}$ is a double coset. Moreover, for $h \in G(s,t)$ and double coset decomposition $w = xhy \in Y_s h Y_t$ we have $\sgn_s(wt) = \sgn(x)\cdot \sgn_s(ht)$.
        \item $\sgn_s(t) \neq 0 \iff \sgn_t(s) \neq 0 \iff \{(s_i, t_i)\}_{i=1,\ldots,m} = \{(i,j) \in D(\lambda) \}$ is the Young diagram of $\lambda$, where $s = s_1\ldots s_m$ and $t = t_1 \ldots t_m$. 
        \item $\sgn_{s}(t) \neq 0 \iff \sgn_{hs}(ht) \neq 0$ for any $h\in S_m$.
        \item $G(s,t)^{-1} = G(t,s)$.
        \item The product $(-1)^{G(s,t)} := \sgn_s(\pi t) \cdot (-1)^{\pi} \cdot \sgn_{t}(\pi^{-1}s)$ does not depend on $\pi \in G(s,t)$, i.e. it is invariant on the double coset.
        \item The product $f(s,t) := (-1)^s(-1)^{G(s,t)}(-1)^t$
        does not depend on a pair $(s,t)$, but only on $\lambda$ and is equal to $(-1)^{\lambda} := (-1)^{t_0}$ for $t_0 = \underbrace{12\ldots \lambda_1}_{1\text{-st block}}\ldots \underbrace{12\ldots\lambda_{\ell}}_{\ell\text{-th block}}$.
    \end{enumerate}
\end{lemma}

\begin{proof}
    We prove the statements one by one.
    
    (a) Due to the weight of the words $s$ and $t$ we have $G(s,t) \neq \varnothing$. Let us fix $h \in G(s,t)$, i.e. $\sgn_s(ht) \neq 0$.
	For an arbitrary permutation $w \in G(s,t)$ with $\sgn_{s}(w t) \neq 0$, there is a unique permutation $x \in Y_{s}$ which permutes elements only in blocks of $s$ so that the words $x\cdot h t = w \cdot t$ are equal. For $y := (xh)^{-1} w$ we have $y \cdot t = t$, and hence $y \in Y_t$. Therefore, the decomposition $w = xhy \in Y_s h Y_t$ holds. 
	
	Conversely, suppose that $w = xhy \in Y_s h Y_t$. Then $y$ leaves $t$ unchanged, $h$ makes $\sgn_s(ht) \neq 0$. Since each transposition of $x$ to the word $hy \cdot t$ changes the sign by $-1$, we have $\sgn_s(wt) = \sgn_s(xht) = (-1)^{x} \cdot \sgn_s(ht) \neq 0$ and $w \in G(s,t)$.
	
	\vspace{0.5em}

    (b) Let us present $s$ and $t$ as a $2 \times m$ table with the first row $s$ and the second row $t$, which we denote by 
    $$(s,t) := \begin{pmatrix}
        s_1 & \ldots & s_m\\
        t_1 & \ldots & t_m
    \end{pmatrix}.
    $$ 
    Then $\sgn_s(t) \neq 0$ if and only if the set of columns of the table $(s,t)$ is equal to the set of cell coordinates of the diagram $\lambda$. 
    
    \vspace{0.5em}

    (c) Swapping rows or columns does not violate this property. 
    
    \vspace{0.5em}

    (d) Let $w^{-1} \in G(s,t)^{-1}$, i.e. $\sgn_s(w^{-1} t) \neq 0$. Then by (b) we have:
    $$
        \sgn_s(w^{-1}t) \neq 0 \iff \sgn_{w^{-1}t}(s) \neq 0 \iff \sgn_{t}(ws) \neq 0,
    $$
    i.e. $w \in G(t,s)$ and all steps are reversible.
    
    \vspace{0.5em}

    (e) By (a) we can write $G(s,t) = Y_s h Y_t$ for an arbitrary element $h \in G(s,t)$. Since Young subgroups are generated by transpositions, this double coset can also be generated with multiplying $h$ by transpositions of $Y_s$ from the left and by transpositions of $Y_t$ from the right. 
    For any transposition $\alpha \in Y_s$, replacing $h \to \alpha h$ we have: $\sgn_s(\alpha h t) = -\sgn(h t)$ (transposition within $s$-block), $\sgn_t((\alpha h)^{-1}s) = \sgn_t(h^{-1} s)$ (transposition of equal letters) and $(-1)^{\alpha h} = -(-1)^{h}$, thus the product is unchanged. 
    Similarly, for a transposition $\beta \in Y_t$ replacing $h \to h \beta$ we have: $\sgn_s(w\beta t) = \sgn_s(wt)$, $\sgn_t((h\beta)^{-1} s) = -\sgn_t(hs)$ and $(-1)^{h\beta} = -(-1)^{h}$. Thus, the product is indeed invariant on the double coset.
    
    \vspace{0.5em}

    (e) Let $s_0 = 1^{\lambda_1}\ldots \ell^{\lambda_\ell}$ and $t_0 = 12\ldots\lambda_1\ldots12\ldots\lambda_\ell$. 
    We will show consequently that $f(s,t) = f(s_0, t_0)$. We will do that in two steps: first, we replace $t \to ht$ so that $\sgn_s(ht) \neq 0$, i.e find a nonzero shift; then we sort the columns of $(s,ht)$ to achieve $(s_0, t_0)$. Diagrammatically we have:
    $$
        \begin{pmatrix}
            s\\
            t
        \end{pmatrix} 
        \xlongrightarrow{t \to ht}
        \begin{pmatrix}
            s\\
            ht
        \end{pmatrix} 
        \xlongrightarrow{\text{ sort columns }}
        \begin{pmatrix}
            s_0\\
            t_0
        \end{pmatrix} = 
        \begin{pmatrix}
            1^{\lambda_1}& \ldots & \ell^{\lambda_\ell}\\
            12\ldots\lambda_1&\ldots& 12\ldots\lambda_\ell
        \end{pmatrix}
    $$
    
    Let $h \in G(s,t)$ be the permutation of minimal length satisfying $\sgn_s(ht) \neq 0$. This permutation does not permute equal letters of $t$ (otherwise the length can be reduced). %hence the relative order of each $s$-block is preserved by $h$. 
    Let us show that $f(s,t) = f(s, ht)$. By (c) we have:
    \begin{align*}
        f(s,t) &= (-1)^s\sgn_s(ht) \cdot (-1)^h \cdot \sgn_t(h^{-1} s) (-1)^t,\\ 
        f(s,ht) &= (-1)^s \sgn_s(ht) \cdot (-1)^{\mathrm{id}} \cdot \sgn_{ht}(s) (-1)^{ht}.
    \end{align*}
    Note that $(-1)^{ht} = (-1)^h (-1)^t$, since each transposition of $h$ swaps only distinct letters of $t$. It remains to show that $\sgn_t(h^{-1} s) = \sgn_{ht}(s)$. Indeed, the table $(t, h^{-1}s)$ is obtained from the table $(ht, s)$ by permuting its columns according to the permutation $h$, which does not permute equal letters of $t$. Hence, the relative order %of permutations in positions 
    of $t$-block permutations in the bottom row does not change in both tables; this implies that $\sgn_t(h^{-1} s) = \sgn_{ht}(s)$ and $f(s,t) = f(s,ht)$.

    Now we may assume that $\sgn_s(t) \neq 0$ (i.e. replace $ht \to t$), then the table $(s,t)$ contains distinct columns forming  the set $\{(i,j)\in D(\lambda)\}$. %Denote $s_0 = 1^{\lambda_1}\ldots\ell^{\lambda_\ell}$ and $t_0 = 12\ldots\lambda_1\ldots12\ldots\lambda_\ell$. 
    By (b) we can reach the table $(s_0,t_0)$ from the table $(s,t)$ by column swaps. Let us consequently apply adjacent transpositions of the form $(i,i+1)$ to the columns of $(s,t)$ to obtain $(s_0, t_0)$ and observe that $f(s,t)$ does not change. By (d), put $\pi = e$ to rewrite
    $$
        f(s,t) = (-1)^s \sgn_s(t) \cdot \sgn_t(s) (-1)^t.
    $$
    Let $\alpha=(i,i+1)$ be the transposition we apply. Let us keep track of the sign changes after applying $\alpha$. If $i$-th and $(i+1)$-th columns are: 
    \begin{itemize}
        \item $\begin{bmatrix}
            xx\\
            yz
        \end{bmatrix}$ then the changes are $\sgn_{\alpha s}(\alpha t) \to -\sgn_s(t)$ and $(-1)^{\alpha t} = -(-1)^t$;
        \item $\begin{bmatrix}
            yz\\
            xx
        \end{bmatrix}$ then the changes are $\sgn_{\alpha t}(\alpha s) \to -\sgn_s(t)$ and $(-1)^{\alpha s} = -(-1)^s$;
        \item $\begin{bmatrix}
            xy\\
            zw
        \end{bmatrix}$ then the changes are $(-1)^{\alpha s} \to -(-1)^s$ and $(-1)^{\alpha t} \to -(-1)^{t}$,
    \end{itemize}
    where $x,y,z,w$ are distinct letters.
    The case where columns are equal is impossible. As we can see in each case $f(s,t)$ does not change, hence the statement follows.
\end{proof}

\end{document}
